\section{Evaluation}
\label{sec:evaluation}

I validated Epistract across six drug discovery domains. Scenarios~1--5 each use a curated corpus of 15--16 PubMed abstracts; Scenario~6 introduces multi-source corpus assembly from PubMed, Google Scholar, and Google Patents via SerpAPI.
Every scenario exercises a distinct facet of the biomedical schema and pipeline, from gene-centric neurogenetics to multi-source competitive intelligence.
Table~\ref{tab:scenarios} summarizes the quantitative results.

\subsection{Six Evaluation Scenarios}

\begin{table}[t]
\centering
\small
\caption{Evaluation scenarios. Each row reports input documents, extracted graph statistics, and user acceptance test (UAT) pass rate. Scenario~6 is the first to use multi-source corpus assembly (PubMed + Google Scholar + Google Patents).}
\label{tab:scenarios}
\begin{tabular}{@{}clrrrrc@{}}
\toprule
\textbf{\#} & \textbf{Domain} & \textbf{Docs} & \textbf{Nodes} & \textbf{Links} & \textbf{Comm.} & \textbf{UATs} \\
\midrule
1 & Neurogenetics       & 15 & 149 & 457 & 6 & 4/4 \\
2 & Oncology            & 16 & 108 & 307 & 4 & 5/5 \\
3 & Rare Disease        & 15 &  94 & 229 & 4 & 3/3 \\
4 & Immuno-Oncology     & 16 & 132 & 361 & 5 & 4/4 \\
5 & CV / Inflammation   & 15 &  94 & 246 & 5 & 3/3 \\
6 & GLP-1 CI$^\dagger$  & 34 & 206 & 630 & 9 & 6/6 \\
\midrule
  & \textbf{Total}      & 111 & 783 & 2{,}230 & 33 & 25/25 \\
\bottomrule
\end{tabular}
\vspace{2pt}
{\footnotesize $^\dagger$Multi-source: 24 PubMed abstracts + 10 patent documents from 5 companies (Novo Nordisk, Eli Lilly, Pfizer, Hanmi Pharmaceutical, Zealand Pharma).}
\end{table}

\textbf{Scenario~1 (Neurogenetics)} targets Alzheimer's disease genetics, centering on the PICALM locus and its network of risk genes.
This corpus is gene-heavy and tests the pipeline's ability to resolve gene--disease associations and group co-implicated loci into coherent communities.

\textbf{Scenario~2 (Oncology)} maps the KRAS G12C inhibitor landscape, including sotorasib, adagrasib, and emerging combination strategies.
It exercises resistance-mechanism relations (\texttt{CONFERS\_RESISTANCE\_TO}) and tests whether the schema captures the interplay between sequence variants, compounds, and clinical trials.

\textbf{Scenario~3 (Rare Disease)} spans three unrelated rare conditions---phenylketonuria (PKU), achondroplasia, and Niemann--Pick type~C---to stress-test schema breadth.
It is the first scenario to exercise \texttt{CONTRAINDICATED\_FOR} and \texttt{GRANTS\_APPROVAL\_FOR}, covering both therapeutic and regulatory relationships.

\textbf{Scenario~4 (Immuno-Oncology)} focuses on checkpoint inhibitor combinations, particularly nivolumab and ipilimumab across melanoma, NSCLC, and renal cell carcinoma indications.
The corpus is protein-heavy and dominated by \texttt{COMBINED\_WITH} and \texttt{BINDS\_TO} relations, testing the pipeline's handling of multi-drug regimens linked to specific clinical trials.

\textbf{Scenario~5 (Cardiovascular / Inflammation)} deliberately includes two unrelated therapeutic areas---hypertrophic cardiomyopathy (HCM, mavacamten) and plaque psoriasis (deucravacitinib)---to test whether community detection correctly separates them into distinct subgraphs without spurious cross-connections.

\textbf{Scenario~6 (GLP-1 Competitive Intelligence)} is the first scenario to use multi-source corpus assembly. The corpus comprises 24 PubMed abstracts and 10 patent documents from five pharmaceutical companies (Novo Nordisk, Eli Lilly, Pfizer, Hanmi Pharmaceutical, Zealand Pharma), assembled via PubMed E-utilities, Google Scholar, and Google Patents through SerpAPI. This scenario tests three novel capabilities: (1)~extraction of molecular identifiers from patent documents---peptide sequences, CAS registry numbers (e.g., tirzepatide: 2023788-19-2), InChIKeys, and chemical formulas; (2)~cross-source entity merging, where semaglutide appearing in 17 documents (4 patents + 13 papers) is correctly deduplicated into a single node with merged attributes; and (3)~competitive landscape community structure, where the graph self-organizes into commercially meaningful clusters (oral delivery technology, triple agonist pipeline, CagriSema combination therapy, non-GLP-1 MASH competitor drugs).

The multi-source assembly workflow mirrors how a human competitive intelligence analyst actually works: searching PubMed for clinical trial abstracts, Google Scholar for broader research papers and patent references, and Google Patents for molecular data and assignee information. This is a design insight---Epistract augments rather than replaces human research workflows, with each source type contributing distinct knowledge: PubMed provides structured clinical data, Scholar enables discovery of papers outside PubMed's index, and patents yield molecular identifiers (sequences, CAS numbers, chemical formulas) unavailable in published abstracts. Publisher paywalls (Lancet, JAMA, NEJM, BMJ) prevented automatic ingestion of several high-impact papers, suggesting a ``Bring Your Own Papers'' model where scientists with institutional access can supplement the corpus.

\subsection{Cross-Domain Schema Coverage}
\label{sec:coverage}

Across all six scenarios, the pipeline exercised 17 of 17 entity types (100\%) and 28 of 30 relation types (93\%).
Each domain naturally activates different schema facets:

\begin{itemize}
  \item \textbf{Scenario~1} is gene-dominant (48 \texttt{GENE} entities), with \texttt{IMPLICATED\_IN} as the most frequent relation type.
  \item \textbf{Scenario~2} is compound-heavy, with \texttt{CONFERS\_RESISTANCE\_TO} linking \texttt{SEQUENCE\_VARIANT} nodes to \texttt{COMPOUND} nodes.
  \item \textbf{Scenario~3} is clinical-heavy and the first to exercise \texttt{CONTRAINDICATED\_FOR} and \texttt{GRANTS\_APPROVAL\_FOR}, covering regulatory-action edges absent from the other corpora.
  \item \textbf{Scenario~4} is protein-heavy, dominated by \texttt{COMBINED\_WITH} and \texttt{BINDS\_TO} relations reflecting checkpoint-receptor biology.
  \item \textbf{Scenario~5} is mechanism-heavy, with \texttt{EVALUATED\_IN} as the dominant relation linking compounds to their pivotal trials.
  \item \textbf{Scenario~6} is compound-dominant (34 \texttt{COMPOUND} entities), with \texttt{INDICATED\_FOR} (64 instances) as the most frequent relation type and \texttt{ACTIVATES} (26 instances) capturing GLP-1 receptor agonism. It is the first scenario to exercise \texttt{DEVELOPS} (company$\to$compound) and \texttt{REGULATES\_EXPRESSION} (semaglutide$\to$GABA in addiction). The \texttt{PEPTIDE\_SEQUENCE} entity type---underexercised in Scenarios~1--5---found its natural home in patent documents, with five validated amino acid sequences.
\end{itemize}

The two unused relation types---\texttt{PHOSPHORYLATES} and \texttt{METABOLIZED\_BY}---represent specialized biochemical relationships (kinase signaling, metabolic clearance) that are rarely discussed at the abstract level.
Full-text extraction would likely activate these relations, but they fall outside the scope of the current evaluation.
Figure~\ref{fig:coverage} presents a heatmap of entity-type and relation-type usage across all six scenarios.

\subsection{Community Detection Quality}
\label{sec:communities}

The Leiden-based community detection step produced 33 communities across the six scenarios, each assigned a biologically meaningful auto-generated label by the summarization agent.
Representative examples illustrate the specificity of the labels:

\begin{itemize}
  \item \emph{``Alzheimer Disease Risk Loci (30 genes)''}---a gene-dominant cluster in Scenario~1 that groups co-implicated GWAS loci around the PICALM hub node.
  \item \emph{``EGFR Inhibitors / Adavosertib / Panitumumab''}---a combination-strategy cluster in Scenario~2 that captures emerging multi-target regimens.
  \item \emph{``TYK2 Allosteric Inhibition''}---a mechanism cluster in Scenario~5 that isolates the deucravacitinib mode-of-action subgraph from the unrelated cardiology cluster.
  \item \emph{``GLP-1(7-37) / SNAC / Semaglutide''}---a drug delivery innovation cluster in Scenario~6 that groups the oral semaglutide formulation with the SNAC absorption enhancer and the native GLP-1 peptide sequence, correctly capturing the technology platform as a distinct community.
  \item \emph{``Denifanstat / Efruxifermin / Lanifibranor''}---a competitive landscape cluster in Scenario~6 that groups non-GLP-1 MASH therapies, emerging from a network meta-analysis paper as a coherent set of competitor compounds.
\end{itemize}

Scenario~5 provided the strongest test of community separation.
Community detection correctly assigned HCM-related nodes (mavacamten, cardiac myosin, obstructive HCM) and psoriasis-related nodes (deucravacitinib, TYK2, plaque psoriasis) to distinct communities with no spurious cross-connections, confirming that the Leiden algorithm~\cite{traag2019louvain} respects topological boundaries even when unrelated literatures are mixed in a single corpus.

Scenario~6 produced 9 communities from 206 nodes---the most complex graph in the evaluation. The communities self-organized into commercially meaningful clusters: MASH therapeutics (survodutide, retatrutide, GCGR agonism), oral delivery technology (SNAC, semaglutide), next-generation competitors (orforglipron, tirzepatide), combination therapy (CagriSema, amylin), emerging CNS indications (addiction, appetite regulation, GABA modulation), and competitor drugs (non-GLP-1 MASH therapies). This demonstrates that the pipeline can produce graphs structured for competitive intelligence analysis, not just scientific review.
Figure~\ref{fig:scenarios} presents multi-panel graph visualizations for all six scenarios.

All 25 user acceptance tests (UATs) passed across the six scenarios.
Each UAT verifies that a specific expected relationship appears in the constructed graph---for example, ``Does \texttt{COMPOUND}(nivolumab) connect to \texttt{CLINICAL\_TRIAL}(CheckMate-067) via \texttt{EVALUATED\_IN}?''
UATs serve as functional regression tests: they confirm that the extraction prompt, schema enforcement, and graph construction pipeline jointly produce the relationships a domain scientist would expect from a given literature set.
